\documentclass[11pt]{article}

\usepackage[margin=1in]{geometry}
\usepackage[hidelinks]{hyperref}
\usepackage{parskip}

\pagestyle{empty}
\setlength{\parindent}{0pt}

\begin{document}

I-Chau Wang\\
Irvine, CA\\
\href{mailto:ichauwang712@gmail.com}{ichauwang712@gmail.com}
\quad\textbar\quad
949-659-7565
\quad\textbar\quad
\href{https://linkedin.com/in/ichauwang}{linkedin.com/in/ichauwang}

\vspace{1.2em}

October 30, 2026

\vspace{1.2em}

Sophinea\\
Mobile App Developer Intern\\
Manassas, VA

\vspace{1.2em}

Dear Sophinea team,

Sophinea has built AI and analytics products for federal and commercial clients since 2018. Halo is the secure platform, and the mobile internship is the layer that puts it on a phone. The work is production screens, API calls, tests, and review. That is close to an app I have already shipped.

NCCUPass is a Flutter app I built for iOS and Android at National Chengchi University. More than 3,000 students use it. One Dart codebase ships both phone platforms. I moved shared logic into 8 services and cut duplicate code by 40\% across 12 screens and 6 features, so a change lived in one place. Every merge request runs formatting, static analysis, and unit tests in GitLab CI.

RailMeet is a second product, written in TypeScript with Next.js and Fastify. The screen calls an API, then draws a ranked meeting place on a map from a PostGIS query. The repository is public on GitHub. The habit underneath both products is the same. A feature is finished when someone can use the result, and the merge request has already been checked. I would bring that same finish line to a Halo screen.

I am in the Master of Computer Science program at UC Irvine through December 2027. I can work part time during the semester on CPT under F-1 status. I would like to talk about the hybrid days in Manassas alongside that semester, and about where the Flutter screens and the API work can help Halo. Thank you for reading.

\vspace{1.2em}

Sincerely,\\
I-Chau Wang

\end{document}
